\part{Hierarchical and Computational Representations}

\section{Filesystems and Shell Environments}

A filesystem state can be modeled as a structure $F=(N,N_{\mathrm{dir}},r,\mathrm{ent},\mathrm{mode})$ in which $N$ is a set of nodes, $N_{\mathrm{dir}}\subseteq N$ the directories, $r\in N_{\mathrm{dir}}$ the root, $\mathrm{ent}:N_{\mathrm{dir}}\times\mathrm{Names}\pfn N$ the directory-entry function, and $\mathrm{mode}:N\to\mathrm{Modes}$ an assignment of permission bits. Symbolic links are nodes that carry a target path, and the entries \texttt{.} and \texttt{..} are ordinary entries of each directory, the latter naming the physical parent. Path resolution is the partial function $\res_F:N_{\mathrm{dir}}\times\mathrm{Paths}\pfn N$ obtained by iterating $\mathrm{ent}$ along the components of a path, expanding symbolic links when they are met. In the special case where the graph of $\mathrm{ent}$, omitting \texttt{.}, \texttt{..}, and links, is a rooted tree, the filesystem is a hierarchy in the sense of Proposition~\ref{prop:treeposet}. Hard links to files make the structure a directed acyclic graph in which a file has several parents, so that the path to a file is no longer unique, and symbolic links add edges that need not respect the tree at all.

The working directory of a shell illustrates the difference between a route and an endpoint recorded in Proposition~\ref{prop:catfromctx}. A POSIX shell maintains a logical working directory as a path string, and the built-in that changes directory behaves differently under the logical option and the physical option when symbolic links are present \cite{posix}. If \texttt{link} in \texttt{/home/u} is a symbolic link to \texttt{/var/data}, then changing into \texttt{link} and then into \texttt{..} returns to \texttt{/home/u} under logical resolution and moves to \texttt{/var} under physical resolution, because the physical parent of \texttt{/var/data} is \texttt{/var}. The parent of a context is therefore not a function of the node alone once links are present: the logical working directory retains an arrow of the path category, and the physical working directory retains only a point of the reachability preorder. A correct model of the context in which a relative path is interpreted must say which of the two the shell uses, and the logical choice is the unfolding of the graph into paths that Section~19 analyzes.

A process adds a second layer that is distinct from the filesystem. A process state includes a working directory, an environment mapping variable names to strings, a file-creation mask, credentials, and open file descriptors. The process tree is generated by spawning: a child created by \texttt{fork} receives a copy of its parent's state at the moment of creation, and a subshell is such a child. Inheritance is therefore a snapshot and not a live reference, since later changes to the parent's environment do not reach an existing child and changes made by the child vanish when it exits. This contrasts with a closure in a lexically scoped language, which captures the location of a binding and so remains coupled to later updates, as Section~6 shows. A third layer is the command interpreter, whose resolution of a command name consults an ordered list: aliases, then shell functions, then built-ins, then the directories named in the search path in order, and the first match wins. This ordered list has the form of a scope chain, but the order is fixed by the interpreter and not by nesting.

Changing the working directory, entering a subshell, and altering file permissions share contextual features, for each changes what some later operation will see, yet they are mathematically distinct. To make this precise, let the configuration of the system be $F\times\mathrm{Proc}^{+}$, a filesystem state together with a nonempty tree of processes with one designated current process. Let $\mathsf{cd}_d$ change the working directory of the current process, defined when the process holds search permission on $d$; let $\mathsf{sub}$ create a child copy of the current process and make it current; let $\mathsf{exit}$ discard the current process and return to its parent; and let $\mathsf{chmod}_{d,m}$ set the mode of $d$ to $m$, defined for the owner.

\begin{proposition}[Distinctness of $\mathsf{cd}$, $\mathsf{sub}$, and $\mathsf{chmod}$; \stN]\label{prop:cdsubchmod}
\begin{enumerate}[label=(\alph*)]
\item The three operations modify different components of the configuration $F\times\mathrm{Proc}^{+}$: $\mathsf{cd}_d$ changes only the working-directory coordinate of the designated process, $\mathsf{sub}$ changes only the process tree and the designation of the current process, and $\mathsf{chmod}_{d,m}$ changes only the mode coordinate of $F$. Their domains of definition nevertheless read other components, since $\mathsf{cd}_d$ consults the modes in $F$ and the credentials of the current process, and $\mathsf{chmod}_{d,m}$ consults the ownership recorded in $F$ against those credentials.
\item They have different persistence. For every finite sequence $w$ of environment-modifying operations and working-directory changes performed between $\mathsf{sub}$ and the matching $\mathsf{exit}$, the composite $\mathsf{exit}\circ w\circ\mathsf{sub}$ restores the parent's coordinates to their former values, whereas $\mathsf{cd}_d$ alters the working directory of the current process and $\mathsf{chmod}_{d,m}$ alters the permissions seen by every process.
\item They do not commute. If an unprivileged process owns a directory $d$ on which it holds search permission, then the sequence $\mathsf{cd}_d$ followed by $\mathsf{chmod}_{d,000}$ is defined, while the sequence $\mathsf{chmod}_{d,000}$ followed by $\mathsf{cd}_d$ is undefined.
\end{enumerate}
\end{proposition}
\begin{proof}
Parts (a) and (b) are read off from the definitions of the operations as maps on $F\times\mathrm{Proc}^{+}$: a child receives a copy of the parent's state, so operations performed on the child leave the parent's coordinates untouched, and discarding the child at $\mathsf{exit}$ restores the parent's private process coordinates. Effects on shared components, such as modifications of $F$ made by the child, persist after the child exits, and (b) is claimed only for the parent's private coordinates. For (c), after $\mathsf{cd}_d$ the process holds $d$ as working directory, and the owner may then clear all permission bits on $d$. In the reverse order, $\mathsf{cd}_d$ requires search permission on the target directory under standard semantics, which the cleared mode denies, so the second step is undefined.
\end{proof}

The three operations therefore cannot be modeled by one kind of arrow in a single category of contexts. Changing directory is a change of position, entering a subshell is a scoped extension of the process component whose effects are discarded on exit, and changing permissions is a change of the authority assignment of the shared component. The table summarizes the profile.

\begin{center}
\begin{tabular}{p{0.2\textwidth}p{0.27\textwidth}p{0.18\textwidth}p{0.22\textwidth}}
\toprule
Operation & Component changed & Persistence & Seen by\\
\midrule
$\mathsf{cd}_d$ & current process, working directory & until next change & this process\\
$\mathsf{sub}$ with environment changes & process tree, child coordinates & until $\mathsf{exit}$ & the child and its descendants\\
$\mathsf{chmod}_{d,m}$ & filesystem, mode of $d$ & until next change & all processes\\
\bottomrule
\end{tabular}
\end{center}

\section{Lexical Scope and Abstract Syntax Trees}

A program is parsed into an abstract syntax tree, and the binding forms of the language, such as \texttt{let}, \texttt{lambda}, function definitions, loops, and comprehensions, introduce scopes within it. Each occurrence of an identifier is a reference to be resolved to a declaration. The tree of syntax is not by itself the structure that resolution consults, because resolution follows binding and not syntactic adjacency, and a more faithful structure is the scope graph \cite{scopegraphs}.

\begin{definition}[Scope graph and resolution policy]\label{def:scopegraph}
A \emph{scope graph} is a directed graph $G=(V,E)$ of scopes in which every vertex $v$ carries a partial binding function $b_v:\mathrm{Names}\pfn\mathrm{Val}$. A \emph{resolution policy} assigns to each pair $(v,n)$ of a scope and a name a finite ordered sequence $\varrho(v,n)=(v_1,\dots,v_k)$ of eligible scopes, each reachable from $v$ along $E$, with $v_1=v$ permitted. The resolution of $n$ at $v$ is
\[
\res(v,n)=b_{v_i}(n)\quad\text{for the least }i\text{ with }n\in\dom b_{v_i},
\]
and is undefined when no such $i$ exists.
\end{definition}

\begin{theorem}[Determinacy of resolution; \stN]\label{thm:determinacy}
For a resolution policy of the form in Definition~\ref{def:scopegraph}, $\res$ is a partial function from pairs $(v,n)$ to values.
\end{theorem}
\begin{proof}
A finite ordered sequence has a unique first position at which a defined binding occurs whenever any such position exists, so the value selected is determined by $(v,n)$.
\end{proof}

The theorem applies to a precisely specified resolution mechanism and not automatically to every language, since languages with dynamic evaluation of names or with ambiguous imports may not supply a finite ordered sequence. Python is a useful test case because its rule is often summarized as local, enclosing, global, built-in, and that summary hides features that distinguish the policy $\varrho$ from simple traversal of ancestors \cite{pyref}. At compile time every name occurring in a function body is classified. A name that is bound anywhere in the body, by assignment, import, definition, or loop variable, is local unless declared otherwise, and this classification covers binding occurrences that follow the use. For a local name, $\varrho(v,n)=(v)$, so that if the binding has not yet been made when the lookup executes, resolution fails and does not continue outward. A name that is used but not bound locally is free, and its eligible scopes are the enclosing function scopes from the innermost outward, then the module scope, then the built-in scope, with enclosing class scopes omitted. The following fragments exhibit both departures from ancestor traversal.
\begin{verbatim}
x = 1
def f():
    print(x)       # UnboundLocalError: x is local to f
    x = 2

x = 'global'
class A:
    x = 'class'
    def m(self):
        return x   # evaluates to 'global', not 'class'
\end{verbatim}
The policy therefore depends on the name through a classification of the entire body, and it skips class scopes. Two further features concern what a resolved binding is. A closure captures the location of a binding and not its value, so that the expression \texttt{[lambda: i for i in range(3)]} produces three functions that all return the final value of \texttt{i}, which is the contrast with the snapshot behavior of process environments noted in Section~5. An import statement binds a module or copies a reference at import time, a snapshot again. Finally, lexical and dynamic scope differ in which parent relation resolution follows. For lexical scope the relevant relation is the static nesting of definitions, and for dynamic scope it is the chain of callers at run time. In the program that defines $x=1$, a function \texttt{f} returning $x$, and a function \texttt{g} that binds $x=2$ and calls \texttt{f}, the call to \texttt{g} yields $1$ under lexical scope and $2$ under dynamic scope. Two distinct parent relations are thus attached to the same set of activation contexts, and conflating them is the source of many incorrect claims that scope is a tree.

\begin{proposition}[Shadowing is non-injective; \stN]\label{prop:shadow}
Let environments be partial maps $E:\mathrm{Names}\pfn\mathrm{Val}$ and let extension be $E\mapsto E[n\mapsto v]$. This operation is not injective. Consequently the observable lookup behavior of an extended environment does not determine the environment that was extended.
\end{proposition}
\begin{proof}
Let $E_1$ and $E_2$ be distinct environments that differ only at the name $n$. Then $E_1[n\mapsto v]=E_2[n\mapsto v]$.
\end{proof}

Shadowing also governs how the information visible at a scope varies with nesting. Define the information available at a scope $v$ as the partial map $K(v):n\mapsto(\text{declaring scope},\,\res(v,n))$, so that $K(v)$ records which declaration a name resolves to.

\begin{proposition}[Visibility under nesting; \stN]\label{prop:visibility}
Let $G$ be a rooted scope tree with the policy $\varrho(v,n)=(v,\mathrm{par}(v),\mathrm{par}^2(v),\dots,\mathrm{root})$ for every name.
\begin{enumerate}[label=(\roman*)]
\item If $w$ is nested in $v$ then $\dom K(v)\subseteq\dom K(w)$: an inner scope resolves at least the names that its container resolves.
\item The graph of $K(v)$ is contained in the graph of $K(w)$ if and only if no name in $\dom K(v)$ is rebound in a scope strictly between $v$ and $w$ or in $w$ itself.
\item Under the Python policy, which omits class scopes from the eligible scopes of nested functions, statement (i) can fail: if a class body $v$ binds a name $n$ that is bound nowhere else, a method $w$ nested in $v$ does not resolve $n$.
\end{enumerate}
\end{proposition}
\begin{proof}
For (i) the chain of ancestors of $v$ is a suffix of the chain of ancestors of $w$, so any scope that supplies a binding at $v$ is eligible at $w$. For (ii), if some name $n\in\dom K(v)$ is rebound at an intermediate scope or at $w$ then $K(w)(n)$ names a different declaring scope, so the pair belonging to $K(v)$ is absent from $K(w)$; conversely if no rebinding occurs then every pair of $K(v)$ persists. For (iii) the class scope is not eligible from the method, which is the second fragment above.
\end{proof}

Visibility is therefore monotone in nesting at the level of domains in the idealized chain policy, monotone at the level of declarations only in the absence of shadowing, and not monotone even at the level of domains in a real language. This is the concrete content behind the abstract Proposition~\ref{prop:nomonotone}.

\begin{proposition}[Scope graphs strictly generalize scope trees; \stN]\label{prop:graphgen}
Every rooted scope tree determines a scope graph, with an edge from each scope to its parent, whose reachability preorder is the ancestor relation. Not every scope graph admits a tree representation that preserves all edges and reachability relations bijectively.
\end{proposition}
\begin{proof}
The first claim is the identification of the parent relation with a graph. For the second, consider the directed acyclic graph with vertices $d,b,c,a$ and edges $d\to b$, $d\to c$, $b\to a$, $c\to a$, in which $d$ has two incomparable parents $b$ and $c$. In a rooted tree the ancestors of any node form a chain by Proposition~\ref{prop:treeposet}, so $b$ and $c$ cannot both be ancestors of $d$. A tree representation therefore either discards one of the edges $d\to b$ and $d\to c$, or duplicates $d$ into two nodes, one below each parent, after which the correspondence on nodes is no longer bijective.
\end{proof}

The diamond is not a contrived case: a class with two base classes in a language with multiple inheritance realizes it, and so does a module reached through two import paths. This supplies a formal reason not to assume that every contextual system is literally a nested tree, a point taken up again in Section~19.

\section{Access Control and Contextual Authority}

The access-matrix model of protection associates to each pair of a principal and an object a set of rights \cite{lampson}. A capability is an unforgeable reference that carries rights and can be passed from one holder to another \cite{dennisvanhorn}. Discretionary access control, as in the permission bits of a Unix filesystem, attaches rights to objects and grants them according to the identity and group of the requester. In a capability discipline there is no ambient authority, so that authority arises only from references actually held \cite{capmyths}. Execution environments inherit part of their authority from their creators, but inheritance is mediated by explicit mechanisms such as dropping privileges, changing identity on execution of a privileged program, and passing descriptors.

The natural question is whether authorization is a functorial assignment to contexts, that is, whether a context-indexed set of admissible operations is monotone with respect to the relevant order, so that $\cA:(C,\le)\to(\mathcal P(\Sigma),\subseteq)$ is a functor. The answer depends entirely on which order is meant, because three distinct orders are attached to the same execution contexts. The first is containment, the nesting of directories or scopes. The second is spawn, the process tree in which each process is created by another. The third is delegation, in which $c\sqsubseteq d$ when the authority of $c$ is obtained by passing a subset of the authority of $d$. Authority is monotone with respect to delegation by construction, and the content of the next proposition is that it need not be monotone with respect to either of the others.

\begin{proposition}[Authority and context orders; \stN]\label{prop:authorders}
There exist a finite set of contexts, three orders on it for containment, spawn, and delegation, and an assignment $\cA$ such that $\cA$ is monotone with respect to delegation and not monotone with respect to containment or spawn. In particular authorization need not decrease with depth in a hierarchy.
\end{proposition}
\begin{proof}
Let the contexts be $u$, a user shell with working directory \texttt{/home}; $l$, a sandboxed child with working directory \texttt{/home/sandbox}; and $s$, a program with elevated authority started by $u$ with working directory \texttt{/home/bin}. Containment gives $l\preceq u$ and $s\preceq u$. Spawn gives $l$ and $s$ as children of $u$. Delegation gives $l\sqsubseteq u$ only, since $s$ holds authority that $u$ did not pass to it. Let $\cA(u)=\{\mathrm{read},\mathrm{write}\}$, $\cA(l)=\{\mathrm{read}\}$, and $\cA(s)=\{\mathrm{read},\mathrm{write},\mathrm{admin}\}$. Then $\cA(l)\subseteq\cA(u)$, consistent with delegation, while $s\preceq u$ and $s$ is a child of $u$ with $\cA(s)\not\subseteq\cA(u)$, so monotonicity fails for containment and for spawn.
\end{proof}

A second distinction concerns path-based and handle-based checks. Reaching a file by name requires search permission on every directory along the path, so the right to reach decreases, in the sense of a conjunction, along the path. After a file has been opened, however, the descriptor is a capability, and later access is checked against the object and the descriptor's rights and not against the ancestors of the original path. The same object thus has authority that is decreasing in depth for name resolution and independent of depth for held references. A limit on any attempt to translate authority structure wholesale is the undecidability of the safety problem for the general access-matrix model, which asks whether a right can leak to a principal that lacks it \cite{hru}; Section~18 returns to what translations can and cannot preserve about prohibitions.

\section{Interactive Worlds as Transition Systems}

An interactive environment is represented as a labeled transition system $\mathcal G=(S,\Sigma,\longrightarrow)$ with $\longrightarrow\subseteq S\times\Sigma\times S$, in which a state records everything that can matter to the next step and an action label records the input or event. Levels, subworlds, menus, portals, and nested interaction states are interpreted as contextual partitions: a surjection $q:S\to C$ assigns each state to the context it occupies, and the \emph{context graph} $\Gamma_q$ has an edge $c\xrightarrow{a}c'$ whenever some transition $s\xrightarrow{a}s'$ has $q(s)=c$ and $q(s')=c'$. Platform games with a map of worlds and levels, games that navigate a graph of cells joined by portals, and strategy games with stacked interface screens are instances of this architecture, and they differ in which kind of nesting they use. Spatial nesting is containment of geometric regions, as of rooms in a floor of a building. Logical nesting is a stack of modes, as of a menu opened over a level. State-dependent accessibility is the presence of an edge only in some states, as of a door that opens only when a key is held.

\begin{proposition}[Context graphs over-approximate; \stE]\label{prop:overapprox}
The map $q$ sends every run of $\mathcal G$ to a run of $\Gamma_q$, so $q$ is a simulation of $\mathcal G$ by $\Gamma_q$. The converse can fail: with initial state $s_0$ and initial context $q(s_0)$, the graph $\Gamma_q$ may contain runs from $q(s_0)$ that no run of $\mathcal G$ from $s_0$ projects to.
\end{proposition}
\begin{proof}
The first claim is the definition of $\Gamma_q$. For the failure, let $S=\{A,B\}\times\{0,1\}$ with the context map $q(r,k)=r$, initial state $s_0=(A,0)$, an action \texttt{take} from $(A,0)$ to $(A,1)$, and an action \texttt{go} from $(A,1)$ to $(B,1)$ only. Then $\Gamma_q$ has the edge $A\xrightarrow{\texttt{go}}B$, so the run consisting of \texttt{go} alone is a run of $\Gamma_q$ from the initial context $A$, but from $s_0=(A,0)$ no run of $\mathcal G$ begins with \texttt{go}.
\end{proof}

State-dependent accessibility is therefore not recoverable from the unguarded context graph, and a faithful representation must either attach guards, which are predicates on $S$, to the edges of $\Gamma_q$, or enlarge the context set to include the relevant part of the state.

Logical nesting raises a different difficulty, because the target of the action that closes a menu is the context from which the menu was opened and not a function of the current context.

\begin{proposition}[Modal return is not an edge of the context graph; \stN]\label{prop:modal}
Let a system have contexts $c_1,\dots,c_n$ with $n\ge2$, from each of which an action \texttt{open} leads to a context $m$, and from $m$ an action \texttt{close} returns to the context from which $m$ was entered. Suppose $c_1$ and $c_2$ are distinguishable, in that some action is enabled at one and not the other. Then no quotient of the system that identifies configurations with the same current context is bisimilar to it. If menus may open menus without bound, the system is not finite-state in the context component.
\end{proposition}
\begin{proof}
In the quotient, $m$ has \texttt{close}-edges to every $c_i$, so the quotient admits the run $c_1,\texttt{open},m,\texttt{close},c_2$, which the system does not. Bisimilarity would require that the configuration $(m,\text{origin }c_1)$ be matched to $m$ and that the quotient move $m\xrightarrow{\texttt{close}}c_2$ be matched by a system move from $(m,\text{origin }c_1)$, whose only \texttt{close}-successor is $c_1$; since $c_1$ and $c_2$ are distinguishable, no matching exists. For unbounded nesting, the return target is the top of a stack of unbounded depth, which is the structure of a pushdown system and not of a finite graph \cite{vpa}.
\end{proof}

When the nesting depth is bounded, the stack can be folded into the context by taking contexts to be bounded sequences, which restores a finite graph at the price of enlarging the context set. This is the first instance of a general fact developed in Section~19, that a graph can be replaced by a hierarchy only by unfolding, and that the unfolding preserves behavior without preserving the identification of repeated contexts.
